\section{实验与结果分析}\label{ux5b9eux9a8cux4e0eux7ed3ux679cux5206ux6790}

第六章已验证系统在任务创建、状态控制、HITL、快照恢复、工具执行和审计追踪等方面具备可运行的工程基础。本章进一步分析
CEBRA-WP
相关策略在批量蛋白质设计工作流中的行为差异。实验目标限定为工作流层的规划、运行时观测、恢复控制和成本控制，不涉及候选蛋白的湿实验功能验证。

本章围绕四个研究问题展开：

\begin{itemize}
\tightlist
\item
  RQ-A1：CEBRA-WP
  的运行时状态、候选重排和动作效用链路是否可执行、可追踪。
\item
  RQ-A2：相比静态单链和固定阈值门控，动态观测是否改变系统行为。
\item
  RQ-A3：Lite belief-state / 轻量信念状态相比 dynamic\_no\_belief\_state
  是否提供额外机制信息。
\item
  RQ-A4：失败和恢复案例是否能由事件日志、快照和候选 metadata 还原。
\end{itemize}

\subsection{实验设计}\label{ux5b9eux9a8cux8bbeux8ba1}

实验采用四组内部消融设计：\texttt{static\_top1}、\texttt{fixed\_threshold\_gate}、\texttt{dynamic\_no\_belief\_state}
和 \texttt{lite\_belief\_state}。四组策略按照运行时介入深度递进：静态
Top-1
只选择静态最高分候选；固定阈值门控组保留阈值门控和局部修补触发；动态观测组保留恢复链路但不启用显式
Lite belief-state / 轻量信念状态；\texttt{lite\_belief\_state}
组启用运行时状态、runtime adjustment 和 action utility。该设计借鉴了 LLM
Agent
中``推理-行动''、候选搜索和失败反馈机制的思想\citep{yao2022react, yao2023tot, shinn2023reflexion}，但实验对象限定为本文系统中的工作流规划与恢复控制。

如图 7-1
所示，实验矩阵由任务集、策略组、评价指标和证据产物四部分组成。任务集提供不同难度、预算和失败压力；策略组用于隔离静态选择、固定阈值门控、动态恢复和
Lite belief-state /
轻量信念状态的影响；评价指标覆盖成功率、首次成功、高代价调用、恢复事件和机制可观测性；证据产物由
run config、event log、snapshot、report 和聚合 CSV 组成。

图 7-1 实验设计框架
插图文件：\texttt{paper/figures/experiment-design-framework.drawio.svg}

图 7-1 的作用是限定本章统计结果的来源。后续表 7-1 至表 7-8 均来自
\texttt{thesis-final-v1-001}
实验矩阵及其配套日志、快照和聚合产物；对应内部证据入口为
\texttt{../thesis-project.dev/docs/experiment/thesis-final-v1-results.md}
和
\texttt{../thesis-project.dev/output/experiment/thesis-final-matrix/thesis-final-v1-001/}。

\textbf{表 7-1 实验矩阵配置表}

{\def\LTcaptype{none} % do not increment counter
\begin{longtable}{@{}
  >{\raggedright\arraybackslash}p{(\linewidth - 2\tabcolsep) * \real{0.5000}}
  >{\raggedright\arraybackslash}p{(\linewidth - 2\tabcolsep) * \real{0.5000}}@{}}
\toprule
\begin{minipage}[b]{\linewidth}\raggedright
项目
\end{minipage} & \begin{minipage}[b]{\linewidth}\raggedright
配置
\end{minipage} \\
\midrule
\endhead
\bottomrule
\endlastfoot
run\_id & \texttt{thesis-final-v1-001} \\
freeze\_id & \texttt{issue209-baseline-freeze-20260326} \\
planner\_provider & \texttt{deepseek-v4-pro} \\
任务覆盖 & 12 个 task\_keys，覆盖 T1 至 T8 共 8 类场景 \\
策略组 & \texttt{static\_top1} / \texttt{fixed\_threshold\_gate} /
\texttt{dynamic\_no\_belief\_state} / \texttt{lite\_belief\_state} \\
repeats &
\texttt{low\_cost\_first=2}，\texttt{standard=2}，\texttt{high\_cost\_sensitive=1} \\
总 runs & 84 runs，每组 21 runs \\
终态结果 & 81 DONE，3 FAILED \\
执行时间 & 约 6 小时 \\
产物路径 &
\texttt{../thesis-project.dev/output/experiment/thesis-final-matrix/thesis-final-v1-001/} \\
\end{longtable}
}

任务集覆盖 8 类场景：简单 de novo
设计、序列评估、多约束稳定性优化、高代价结构预测、可修复参数失败、远程服务降级、结构性重规划和安全探测。涉及的公开结构包括
Trp-cage、Villin HP35、GB1、Ubiquitin、Top7 和 de novo
oligomer。任务定义与分层依据来自
\texttt{../thesis-project.dev/docs/experiment/final-thesis-experiment-design.md}
和
\texttt{../thesis-project.dev/docs/experiment/thesis-final-v1-results.md}。该任务集使实验同时包含低成本成功路径、中等压力路径和高代价结构预测路径。

\subsection{指标定义}\label{ux6307ux6807ux5b9aux4e49}

本章使用五类指标。任务完成指标包括
\texttt{success\_rate}、\texttt{first\_pass\_success\_rate}、\texttt{schema\_valid\_rate}
和 \texttt{executable\_plan\_rate}。恢复行为指标包括
\texttt{patch\_events\_mean}、\texttt{replan\_events\_mean} 和
\texttt{suffix\_replan\_events\_mean}。成本指标包括
\texttt{high\_cost\_call\_mean}、\texttt{high\_cost\_call\_total} 和
\texttt{duration\_ms\_mean}，其中高代价调用在本实验中主要对应结构预测与结构精修。机制可观测性指标包括
\texttt{runtime\_state\_observable\_rate}、\texttt{runtime\_state\_summary}、\texttt{budget\_pressure\ source}
和 \texttt{action\_utility\_source}。增量指标使用相邻策略组的 paired
delta。

局部修补/重规划事件来自 event log，表示实际执行过的恢复动作；action
utility 来自候选 metadata，表示算法对
\texttt{continue}、\texttt{patch\_local}、\texttt{suffix\_replan} 和
\texttt{stop} 等动作的评分。两者含义不同，后文表格分别列示。

\subsection{四组消融主结果}\label{ux56dbux7ec4ux6d88ux878dux4e3bux7ed3ux679c}

表 7-2 汇总四组策略的核心指标。每组 21 runs，static\_top1 全部
DONE，其余三组各 1 个 FAILED。

\textbf{表 7-2 四组消融主实验结果}

{\def\LTcaptype{none} % do not increment counter
\begin{longtable}{@{}
  >{\raggedright\arraybackslash}p{(\linewidth - 8\tabcolsep) * \real{0.1579}}
  >{\raggedleft\arraybackslash}p{(\linewidth - 8\tabcolsep) * \real{0.2105}}
  >{\raggedleft\arraybackslash}p{(\linewidth - 8\tabcolsep) * \real{0.2105}}
  >{\raggedleft\arraybackslash}p{(\linewidth - 8\tabcolsep) * \real{0.2105}}
  >{\raggedleft\arraybackslash}p{(\linewidth - 8\tabcolsep) * \real{0.2105}}@{}}
\toprule
\begin{minipage}[b]{\linewidth}\raggedright
指标
\end{minipage} & \begin{minipage}[b]{\linewidth}\raggedleft
static\_top1
\end{minipage} & \begin{minipage}[b]{\linewidth}\raggedleft
fixed\_threshold\_gate
\end{minipage} & \begin{minipage}[b]{\linewidth}\raggedleft
dynamic\_no\_belief\_state
\end{minipage} & \begin{minipage}[b]{\linewidth}\raggedleft
lite\_belief\_state
\end{minipage} \\
\midrule
\endhead
\bottomrule
\endlastfoot
runs & 21 & 21 & 21 & 21 \\
DONE & 21 & 20 & 20 & 20 \\
FAILED & 0 & 1 & 1 & 1 \\
success\_rate & 1.0000 & 0.9524 & 0.9524 & 0.9524 \\
first\_pass\_success\_rate & 1.0000 & 0.9048 & 0.9524 & 0.9524 \\
schema\_valid\_rate & 1.0000 & 1.0000 & 0.9524 & 1.0000 \\
executable\_plan\_rate & 1.0000 & 0.9524 & 1.0000 & 1.0000 \\
waiting\_chain\_complete\_rate & 1.0000 & 1.0000 & 1.0000 & 1.0000 \\
failure\_traceable\_rate & 1.0000 & 1.0000 & 1.0000 & 1.0000 \\
runtime\_state\_observable\_rate & 0.0000 & 0.0000 & 0.0000 & 1.0000 \\
patch\_events\_mean & 0.0000 & 0.2857 & 0.0000 & 0.0000 \\
replan\_events\_mean & 0.0000 & 0.0000 & 0.0000 & 0.0000 \\
suffix\_replan\_events\_mean & 0.0000 & 0.0000 & 0.0000 & 0.0000 \\
high\_cost\_call\_mean & 1.0000 & 1.3333 & 0.9524 & 0.9524 \\
high\_cost\_call\_total & 21 & 28 & 20 & 20 \\
duration\_ms\_mean & 241,905 & 300,238 & 226,571 & 272,095 \\
action\_continue\_mean & 1.0952 & 1.5238 & 0.9524 & 0.9524 \\
action\_patch\_local\_mean & 0.0000 & 0.2857 & 0.0000 & 0.0000 \\
\end{longtable}
}

表 7-2 显示出三个主要现象。static\_top1
在该矩阵中成功率最高，提示任务集中的多数初始候选已经具备较强可执行性。fixed\_threshold\_gate
是唯一触发真实局部修补的组，平均局部修补次数为
0.2857，同时高代价调用总数升至 28。lite\_belief\_state 是唯一
runtime\_state\_observable\_rate 为 1.0000 的组，说明其 Lite
belief-state / 轻量信念状态链路在所有 run 中均产生可追踪输出。

\subsection{分层结果与机制增量}\label{ux5206ux5c42ux7ed3ux679cux4e0eux673aux5236ux589eux91cf}

表 7-3 按任务难度和预算层级汇总成功率。3 个 FAILED run 全部位于
medium/standard 层。

\textbf{表 7-3 按任务难度与预算分层的成功率}

{\def\LTcaptype{none} % do not increment counter
\begin{longtable}{@{}lrrrr@{}}
\toprule
分层 & runs & DONE & FAILED & success\_rate \\
\midrule
\endhead
\bottomrule
\endlastfoot
easy & 32 & 32 & 0 & 1.0000 \\
medium & 40 & 37 & 3 & 0.9250 \\
hard & 12 & 12 & 0 & 1.0000 \\
low\_cost\_first & 32 & 32 & 0 & 1.0000 \\
standard & 40 & 37 & 3 & 0.9250 \\
high\_cost\_sensitive & 12 & 12 & 0 & 1.0000 \\
\end{longtable}
}

分层结果说明，失败主要来自中等难度、标准预算任务。具体任务中，\texttt{t2\_ubiquitin\_sequence\_eval}
贡献 2 个 FAILED，\texttt{t3\_gb1\_stability\_optimization} 贡献 1 个
FAILED。hard 层样本量较小，该行仅作为本矩阵中的观察事实。

表 7-4 给出相邻策略组的 paired delta。delta
为后一个策略组减前一个策略组。

\textbf{表 7-4 机制增量配对对比}

{\def\LTcaptype{none} % do not increment counter
\begin{longtable}{@{}
  >{\raggedright\arraybackslash}p{(\linewidth - 8\tabcolsep) * \real{0.1765}}
  >{\raggedright\arraybackslash}p{(\linewidth - 8\tabcolsep) * \real{0.1765}}
  >{\raggedleft\arraybackslash}p{(\linewidth - 8\tabcolsep) * \real{0.2353}}
  >{\raggedleft\arraybackslash}p{(\linewidth - 8\tabcolsep) * \real{0.2353}}
  >{\raggedright\arraybackslash}p{(\linewidth - 8\tabcolsep) * \real{0.1765}}@{}}
\toprule
\begin{minipage}[b]{\linewidth}\raggedright
对比
\end{minipage} & \begin{minipage}[b]{\linewidth}\raggedright
指标
\end{minipage} & \begin{minipage}[b]{\linewidth}\raggedleft
delta
\end{minipage} & \begin{minipage}[b]{\linewidth}\raggedleft
样本数
\end{minipage} & \begin{minipage}[b]{\linewidth}\raggedright
结果含义
\end{minipage} \\
\midrule
\endhead
\bottomrule
\endlastfoot
static→fixed & success\_rate & -0.0476 & 21 & fixed 组多 1 个失败。 \\
static→fixed & first\_pass\_success\_rate & -0.0952 & 21 & fixed
组首次成功率下降。 \\
static→fixed & patch\_event\_count & +0.2857 & 21 &
\texttt{fixed\_threshold\_gate} 组触发真实局部修补。 \\
static→fixed & high\_cost\_call\_count & +0.3333 & 21 & fixed
组高代价调用增加。 \\
fixed→dynamic & high\_cost\_call\_count & -0.3810 & 21 & dynamic
组高代价调用低于 fixed。 \\
fixed→dynamic & duration\_ms & -73,667 & 21 & dynamic 组平均耗时低于
fixed。 \\
dynamic→lite & schema\_valid\_rate & +0.0476 & 21 & lite 组 schema valid
恢复到 1.0000。 \\
dynamic→lite & high\_cost\_call\_count & 0.0000 & 21 &
两组高代价调用均值相同。 \\
dynamic→lite & duration\_ms & +45,524 & 21 & lite 组平均耗时高于
dynamic。 \\
\end{longtable}
}

表 7-4 表明，fixed\_threshold\_gate
引入了可观测的局部修补行为，同时带来额外高代价调用。dynamic\_no\_belief\_state
相比 fixed\_threshold\_gate 降低了高代价调用和耗时。lite\_belief\_state
相比 dynamic\_no\_belief\_state 的主要差异，则体现在 schema
valid、RuntimeState 和 action utility 等机制信息上。

\subsection{成本与恢复行为分析}\label{ux6210ux672cux4e0eux6062ux590dux884cux4e3aux5206ux6790}

表 7-5 单独列出高代价调用和运行时间。fixed\_threshold\_gate
的高代价调用总数为 28，其中结构映射 27 次，结构精修 1
次；dynamic\_no\_belief\_state 和 lite\_belief\_state 均为 20 次。

\textbf{表 7-5 高代价调用与运行时间对比}

{\def\LTcaptype{none} % do not increment counter
\begin{longtable}{@{}
  >{\raggedright\arraybackslash}p{(\linewidth - 10\tabcolsep) * \real{0.1304}}
  >{\raggedleft\arraybackslash}p{(\linewidth - 10\tabcolsep) * \real{0.1739}}
  >{\raggedleft\arraybackslash}p{(\linewidth - 10\tabcolsep) * \real{0.1739}}
  >{\raggedleft\arraybackslash}p{(\linewidth - 10\tabcolsep) * \real{0.1739}}
  >{\raggedleft\arraybackslash}p{(\linewidth - 10\tabcolsep) * \real{0.1739}}
  >{\raggedleft\arraybackslash}p{(\linewidth - 10\tabcolsep) * \real{0.1739}}@{}}
\toprule
\begin{minipage}[b]{\linewidth}\raggedright
组
\end{minipage} & \begin{minipage}[b]{\linewidth}\raggedleft
high\_cost\_total
\end{minipage} & \begin{minipage}[b]{\linewidth}\raggedleft
high\_cost\_mean
\end{minipage} & \begin{minipage}[b]{\linewidth}\raggedleft
结构映射
\end{minipage} & \begin{minipage}[b]{\linewidth}\raggedleft
结构精修
\end{minipage} & \begin{minipage}[b]{\linewidth}\raggedleft
平均耗时
\end{minipage} \\
\midrule
\endhead
\bottomrule
\endlastfoot
static\_top1 & 21 & 1.0000 & 21 & 0 & 241,905 ms \\
fixed\_threshold\_gate & 28 & 1.3333 & 27 & 1 & 300,238 ms \\
dynamic\_no\_belief\_state & 20 & 0.9524 & 20 & 0 & 226,571 ms \\
lite\_belief\_state & 20 & 0.9524 & 20 & 0 & 272,095 ms \\
\end{longtable}
}

由表 7-5 可计算得出，与 fixed\_threshold\_gate
相比，dynamic\_no\_belief\_state 和 lite\_belief\_state
的高代价调用总数从 28 降至 20，降幅为 28.6\%。与 static\_top1
相比，两组高代价调用总数从 21 降至 20，差异来自各自 1 个失败 run
在高代价结构预测前终止。运行时间上，dynamic\_no\_belief\_state
最短，fixed\_threshold\_gate 最长，lite\_belief\_state 介于二者之间。

恢复事件方面，四组均未产生真实重规划或后缀重规划；真实局部修补仅出现在
fixed\_threshold\_gate 组，总数 6 次，均为 tool-level
patch。该结果说明，本次矩阵对局部修补机制产生了真实压力，但对重规划机制的矩阵级触发不足。第六章中的
focused tests 已覆盖后缀重规划和 \texttt{terminal\_stop}
的可达性，因此本章的批量实验结论以局部修补和高代价调用为主。

\subsection{Lite belief-state
机制可观测性}\label{lite-belief-state-ux673aux5236ux53efux89c2ux6d4bux6027}

表 7-6 对比 Lite belief-state 与其余三组的机制可观测性。该表对应 RQ-A1
和 RQ-A3。

\textbf{表 7-6 Lite belief-state 机制可观测性对比}

{\def\LTcaptype{none} % do not increment counter
\begin{longtable}{@{}
  >{\raggedright\arraybackslash}p{(\linewidth - 4\tabcolsep) * \real{0.3333}}
  >{\raggedright\arraybackslash}p{(\linewidth - 4\tabcolsep) * \real{0.3333}}
  >{\raggedright\arraybackslash}p{(\linewidth - 4\tabcolsep) * \real{0.3333}}@{}}
\toprule
\begin{minipage}[b]{\linewidth}\raggedright
特征
\end{minipage} & \begin{minipage}[b]{\linewidth}\raggedright
static\_top1 / fixed\_threshold\_gate / dynamic\_no\_belief\_state
\end{minipage} & \begin{minipage}[b]{\linewidth}\raggedright
lite\_belief\_state
\end{minipage} \\
\midrule
\endhead
\bottomrule
\endlastfoot
runtime\_state\_summary & 无有效运行时状态摘要 & 21/21 runs 产生
RuntimeState \\
runtime\_state\_observable\_rate & 0.0000 & 1.0000 \\
belief-state 核心字段 & 未观测 &
\texttt{p\_success}、\texttt{p\_structural\_failure}、\texttt{recovery\_margin}、\texttt{expected\_remaining\_cost}、\texttt{evidence\_sufficiency} \\
budget\_pressure source & \texttt{default} & \texttt{observed} \\
action\_utility\_source & \texttt{missing} & \texttt{computed} \\
action\_utilities & 空对象或不可用 &
\texttt{continue}、\texttt{patch\_local}、\texttt{suffix\_replan}、\texttt{stop}
均有 utility \\
high\_cost\_call\_mean & static=1.0000，fixed=1.3333，dynamic=0.9524 &
0.9524 \\
\end{longtable}
}

Lite belief-state / 轻量信念状态的核心证据是 21/21 runs 均产生有效
RuntimeState，且 budget pressure
来自运行时观测而非默认回退。该结果说明，CEBRA-WP
的运行时状态更新、预算压力估计和动作效用计算均已在矩阵实验中执行。与
dynamic\_no\_belief\_state 相比，lite\_belief\_state 的 success\_rate 和
high\_cost\_call\_mean 相同，其增量主要体现为完整的 RuntimeState
和动作效用记录。

\subsection{典型失败案例}\label{ux5178ux578bux5931ux8d25ux6848ux4f8b}

如图 7-2 所示，fixed\_threshold\_gate 和 lite\_belief\_state
的恢复控制路径不同。fixed\_threshold\_gate
主要体现为运行时门控后的局部修补；lite\_belief\_state
则在候选评估阶段产生 RuntimeState、runtime adjustment 和 action
utility，并据此影响候选排序。

图 7-2 固定修补与 Lite belief-state 重排序的恢复路径对比
插图文件：\texttt{paper/figures/recovery-path-comparison-timeline.drawio.svg}

图 7-2 用于解释表 7-7 中的失败案例。fixed\_threshold\_gate
的失败体现为局部修补循环耗尽；lite\_belief\_state
的失败体现为运行时状态可观测但仍未打破循环；dynamic\_no\_belief\_state
的失败体现为候选 I/O 闭包校验失败。

\textbf{表 7-7 失败案例归因表}

{\def\LTcaptype{none} % do not increment counter
\begin{longtable}{@{}
  >{\raggedright\arraybackslash}p{(\linewidth - 10\tabcolsep) * \real{0.1579}}
  >{\raggedright\arraybackslash}p{(\linewidth - 10\tabcolsep) * \real{0.1579}}
  >{\raggedleft\arraybackslash}p{(\linewidth - 10\tabcolsep) * \real{0.2105}}
  >{\raggedright\arraybackslash}p{(\linewidth - 10\tabcolsep) * \real{0.1579}}
  >{\raggedright\arraybackslash}p{(\linewidth - 10\tabcolsep) * \real{0.1579}}
  >{\raggedright\arraybackslash}p{(\linewidth - 10\tabcolsep) * \real{0.1579}}@{}}
\toprule
\begin{minipage}[b]{\linewidth}\raggedright
策略组
\end{minipage} & \begin{minipage}[b]{\linewidth}\raggedright
任务
\end{minipage} & \begin{minipage}[b]{\linewidth}\raggedleft
repeat
\end{minipage} & \begin{minipage}[b]{\linewidth}\raggedright
错误类型
\end{minipage} & \begin{minipage}[b]{\linewidth}\raggedright
根因摘要
\end{minipage} & \begin{minipage}[b]{\linewidth}\raggedright
证据意义
\end{minipage} \\
\midrule
\endhead
\bottomrule
\endlastfoot
fixed\_threshold\_gate & \texttt{t2\_ubiquitin\_sequence\_eval} & r02 &
auto decision loop exhausted & 固定阈值门控触发 WAITING\_PATCH
循环，局部修补后仍触发相同门控条件。 &
说明固定阈值门控可以触发恢复，也可能引入额外循环和高代价调用。 \\
lite\_belief\_state & \texttt{t2\_ubiquitin\_sequence\_eval} & r01 &
auto decision loop exhausted & 多次 WAITING\_PATCH 循环中
budget\_pressure 升至 1.5，runtime adjustment 已产生但未打破循环。 &
说明 belief-state 链路可观测，且该机制仍受候选质量和循环控制限制。 \\
dynamic\_no\_belief\_state & \texttt{t3\_gb1\_stability\_optimization} &
r01 & \texttt{CANDIDATE\_IO\_CLOSURE\_BROKEN} & 候选 step
输出字段无法被下游引用，候选验证在执行前失败。 &
说明候选可执行性校验能够阻断无效 plan。 \\
\end{longtable}
}

三个失败案例均有 event log 和
snapshot，可作为可追溯的边界样本。fixed\_threshold\_gate 的 t2
案例解释了表 7-5 中高代价调用增加的来源；lite\_belief\_state 的 t2
案例展示了 RuntimeState 的动态演化；dynamic\_no\_belief\_state 的 t3
案例对应表 7-2 中 schema\_valid\_rate=0.9524。

\subsection{证据产物与本章结论}\label{ux8bc1ux636eux4ea7ux7269ux4e0eux672cux7ae0ux7ed3ux8bba}

表 7-8 汇总本章使用的实验产物。所有核心结论均可回溯到矩阵报告、CSV
聚合、run 级结果、事件日志或快照。

\textbf{表 7-8 实验结果证据产物索引}

{\def\LTcaptype{none} % do not increment counter
\begin{longtable}{@{}
  >{\raggedright\arraybackslash}p{(\linewidth - 6\tabcolsep) * \real{0.2500}}
  >{\raggedright\arraybackslash}p{(\linewidth - 6\tabcolsep) * \real{0.2500}}
  >{\raggedright\arraybackslash}p{(\linewidth - 6\tabcolsep) * \real{0.2500}}
  >{\raggedright\arraybackslash}p{(\linewidth - 6\tabcolsep) * \real{0.2500}}@{}}
\toprule
\begin{minipage}[b]{\linewidth}\raggedright
证据产物
\end{minipage} & \begin{minipage}[b]{\linewidth}\raggedright
路径/来源
\end{minipage} & \begin{minipage}[b]{\linewidth}\raggedright
支撑内容
\end{minipage} & \begin{minipage}[b]{\linewidth}\raggedright
使用位置
\end{minipage} \\
\midrule
\endhead
\bottomrule
\endlastfoot
实验结果报告 &
\texttt{../thesis-project.dev/docs/experiment/thesis-final-v1-results.md}
& 84-run 概况、主结果、分层结果、失败案例 & 全章 \\
实验设计文档 &
\texttt{../thesis-project.dev/docs/experiment/final-thesis-experiment-design.md}
& 研究问题、任务类、指标体系 & 7.1、7.2 \\
分组映射文档 &
\texttt{../thesis-project.dev/docs/experiment/algorithm-group-paper-mapping.md}
& 代码 policy mode 与论文组名对应关系 & 7.1 \\
主指标汇总 & \texttt{matrix\_metrics\_summary.csv} &
四组主指标、成功率、恢复、耗时和高代价调用 & 7.3、7.6 \\
机制增量表 & \texttt{mechanism\_increment\_deltas.csv} & 相邻策略组
paired delta & 7.4 \\
高代价调用表 & \texttt{high\_cost\_breakdown.csv} &
高代价调用总数和规则命中 & 7.5 \\
恢复事件表 &
\texttt{patch\_replan\_breakdown.csv}、\texttt{action\_distribution.csv}
& patch/replan/action 分布 & 7.5 \\
run 级结果 & \texttt{run\_level\_results.jsonl} &
失败案例、RuntimeState、tool usage、artifact linkage & 7.6、7.7 \\
事件日志与快照 &
\texttt{data/logs/thesis-final-v1-001\_*.jsonl}、\texttt{data/snapshots/thesis-final-v1-001\_*.jsonl}
& 决策链、恢复链、运行时状态和失败追踪 & 7.7 \\
\end{longtable}
}

本章实验得到四点结论。在当前实验矩阵中，CEBRA-WP
的机制链路可执行、可追踪，lite\_belief\_state 组 21/21 runs 产生
RuntimeState，runtime\_state\_observable\_rate 为
1.0000。fixed\_threshold\_gate 触发 6 次真实局部修补，高代价调用总数达到
28，说明固定阈值门控能够暴露恢复需求，同时也会带来额外执行成本。dynamic\_no\_belief\_state
与 lite\_belief\_state 的 high\_cost\_call\_mean 均为 0.9524，低于
fixed\_threshold\_gate 的 1.3333；lite 的主要增量体现在 Lite
belief-state / 轻量信念状态、预算压力和动作效用的可观测性。3 个 FAILED
run 均有可追溯事件链，失败集中在 medium/standard
层，主要反映候选验证、局部修补循环和运行时控制的边界条件。

在当前 \texttt{thesis-final-v1-001} 设置下，实验支持``CEBRA-WP
机制已实现且可观测''\,``固定阈值门控恢复存在额外成本''\,``Lite
belief-state /
轻量信念状态提供运行时决策解释信息''等结论。成功率方面，static\_top1 为
1.0000，其余三组为 0.9524，最终成功率提升并非本章的主要实验结论。

第八章将结合上述实验结论与边界，总结本文贡献并讨论后续改进方向。
